\subsection{Six deletions for a path and a disjoint edge}
The path configuration permits a further exclusion when triangles are also
absent. The additional restriction concerns how distinct inserted values
can occur on opposite sides of a quadratic product.

\begin{proposition}[An optimal frame for a missing path and edge]\label{prop:pathsix}
Let $G$ have girth at least five, and let a positive variable-input binary
circuit compute $P\in\Sup(G)$ with nonempty quadratic support $H$.
Suppose its four missing edges form $P_4\sqcup K_2$, the cover gain is
three, and the extracted quadratic circuit is optimal. Then six deletions
are impossible. Consequently every circuit with this missing-edge shape,
cover gain three and an optimal quadratic frame satisfies the full star
lower bound.
\end{proposition}
\begin{proof}
Write the missing edges as $ab,bc,cd,ef$ and call their graph $J$.
The assigned minimum $H$-cover $Z$ avoids all six endpoints, as in
Proposition~\ref{prop:pathfive}. On $T=V(G)\setminus Z$, the graph is
therefore $J$ with possible isolated vertices. Full-carrier erasure loses
no private gate: a loss would leave at most five deletions with the same
optimal frame, already excluded by Proposition~\ref{prop:pathfive} and
the preceding four-witness deletion bound. This assertion concerns the
original circuit; restricted private values may be zero, redundant or
equal.

We use three restrictions from the quadratic frame. Every nonzero
nonconstant individual direct context on $T$ is not contained in a
minimum cover of $J$. Indeed, such a context contains all outside
neighbors of the literal assigned center it perturbs; replacing that
center by a minimum $J$-cover containing the context would cover $G$
with $\tau(H)+2$ vertices. A union of permitted contexts obeys the same
condition. Second, an old linear boundary clause and an individual
linear context avoid the pairs $\{a,c\}$ and $\{b,d\}$. A common old
neighbor for either pair would complete a four-cycle with its missing
two-edge path. The context assertion uses the positive old product
$AL\subseteq P_2$, where $A$ is the cut's nonzero old low part.
Third, triangle-freeness also prohibits adjacent vertices of $J$ in an
old clause or individual context: their common old neighbor would form
a triangle. We use this third restriction only in the case of three
distinct visible highs. Neither pair restriction is asserted for a
union of contexts from different cuts or for the output of a private
addition.

Identify all isolated vertices of $T$ with one variable $z$. The target
remains $J$, and the restrictions on its six genuine endpoints persist.
Zero and constant-one affine contexts are allowed throughout. With at
most two private products, the unrestricted two-product class in
Proposition~\ref{prop:pathfive} already gives a contradiction, regardless
of the number of private additions. With one deleted addition,
Theorem~\ref{thm:purecut} gives $D\ge 2^3-1=7$. At least one addition is
deleted. The remaining product/addition counts are $(4,2)$ and $(3,3)$.

\emph{Four products and two additions.}
If only one addition is visible, its high value is computed by a
five-node Boolean DAG with four AND gates and one OR gate over old
linear clauses. Every private node is ancestral to that high, because
it has no other exit to the selected quadratic subgraph. Its one
affine context obeys both the minimum-cover and missing-path
restrictions. This is the sixth finite class below.

If both additions are visible and select the same high $g$, no private
OR remains. The four-product multiplication DAG has at most five
distinct old boundary factors. Its merged context gives at most six
CNF clauses. The edge function of $J$ has six minimal vertex covers,
all of size three. Every positive CNF representing this function must
contain each of those covers as a clause: at its complementary maximal
independent assignment, some false clause is a cover contained in the
minimal cover, and hence equals it. A representation with at most six
clauses must therefore have exactly these clauses. A nonconstant
merged context cannot be one of them by the minimum-cover restriction;
a constant-one context leaves at most five clauses. Both alternatives
are impossible.

If the two highs are distinct, they are private products. Expose each
as soon as it is computed, which only enlarges its actual availability.
A private operand is an old linear clause joined with earlier exposed
highs, or one earlier private product alone. An unexposed product cannot
mix with another argument without a private addition, and none is
available. Each individual context obeys the two restrictions. Only
the two exposed indices may occur in direct output terms or pair
terms. This is the seventh finite class.

\emph{Three products and one visible cut.}
The high computation is a five-node DAG with three AND gates and two
private OR gates. Each input is an old clause or an earlier node, with
arbitrary sharing. Old clauses satisfy the missing-path restriction;
private OR outputs are unrestricted. The unique affine context
satisfies both restrictions. This is the first finite class.

\emph{Three products and two visible cuts.}
There is one private OR node $q$. If the highs coincide, all four
private nodes precede their common high, whose computation has three
ANDs and one OR. Its merged context obeys the minimum-cover condition
but need not avoid a missing-path pair. Diagonal terms are absorbed
into a constant-one context. This is the third finite class.

If the highs are distinct, each is a private product or $q$. An old
linear value or a degree-one visible cut is too small to be a high
operand. A degree-two visible cut is a protected carrier and is
unavailable by erasure. Thus the three products and $q$ form four
private nodes, at most two of which are designated highs. Each input
is an old clause joined with earlier designated highs, or an earlier
private node alone. The OR node is explicit, including when it joins
two old clauses; its output is not subjected to an old-clause pair
restriction. Each individual output context obeys both restrictions.
This is the second finite class.

\emph{Three products and three visible cuts, with repeated highs.}
No private addition remains, and every visible high is a private
product. If all highs coincide, at most four old factors and one
merged context give five clauses, fewer than the required six. If
there are exactly two distinct highs, one is used by two cuts and one
by a single cut. The former's merged context obeys only the
minimum-cover condition; the latter's individual context also avoids
the missing-path pairs. There is at most one mixed pair term, and
diagonal terms are absorbed into constant contexts. Unselected
products remain available alone but cannot mix with old arguments.
The fifth finite class permits at most two exposed output indices and
allows at most one context to violate the missing-path restriction.

\emph{Three products and three distinct visible highs.}
Here each private product is the high operand of one visible cut.
Let $S_i$ consist of the genuine missing endpoints $u$ on whose
singleton assignment the restricted high $g_i$ is true. The identified
isolate $z$ is excluded from $S_i$.

Call a cut center-affecting if its independent degree-one symbol
reaches the literal center factor of a quadratic star. Its old low
part is that literal center $h\in Z$, and its restricted direct
context contains $N_H(h)\cap T$, a nonempty set because $Z$ is minimum.
If $g_i$ were true on a singleton $u$, safety of its direct contribution
would put that context inside $N_J(u)$. Every neighborhood in
$P_4\sqcup K_2$ is contained in a minimum cover of $J$: the two internal
path neighborhoods are $\{a,c\}$ and $\{b,d\}$, and the other
neighborhoods are singletons extendable to such a cover. This
contradicts the context condition. Therefore
\[
 \text{a center-affecting cut has }S_i=\varnothing.             \tag{*}
\]

If the independent-symbol pair coefficient $\delta_{ij}$ is positive,
the two symbols enter opposite factors of a retained quadratic
product. One factor is its literal center, so at least one cut is
center-affecting. By $(\ast)$, $S_i$ and $S_j$ cannot both be nonempty.
A degree-two cut has scalar context and cannot create a quadratic
pair term.

There is also a common-neighbor restriction. Fix $j$ and put
$U_j=\bigcup_{i:\delta_{ij}>0}S_i$. If cut $j$ is never
center-affecting, every adjacent cut $i$ is center-affecting at the
corresponding product, and $U_j$ is empty by $(\ast)$. Otherwise the
old low part of cut $j$ is one fixed literal $h_j\in Z$ at every use.
Each pair term $g_i g_j$ then has a positive contribution $g_i h_j$
obtained by taking the old low value on the $j$ side. Safety on
$\{u,h_j\}$, for $u\in S_i$, forces the edge $uh_j$. It is an old edge
because every missing endpoint lies outside $Z$. Hence $U_j$ is
contained in one old neighborhood. It contains neither adjacent
vertices of $J$ nor the outer pair of a missing two-edge path, by
triangle- and four-cycle-freeness.

These two singleton restrictions require the distinct-cut hypothesis
just used. The fourth finite class imposes them only in the present
case. It otherwise allows all three products in later input
disjunctions, all individual affine contexts and all pair terms.
Old clauses and contexts avoid both the missing-path pairs and the
adjacent pairs, and contexts obey the minimum-cover condition.

\emph{Finite verification.}
In each class, Boolean selectors choose old arguments, permitted
earlier high values, contexts and pair terms. Exact bidirectional
Tseitin equivalences evaluate every selected argument and gate.
Truth on the four edge assignments with $z=0$ and falsity on the six
maximal independent assignments with $z=1$ are necessary and
sufficient for equality with the monotone target. The fourth class
also evaluates the six genuine singleton assignments to express its
two conditional restrictions.

For a one-cut DAG, one-hot modes select an old clause or a prior
node, and operation bits impose the exact AND and OR counts.
The class with one private OR and two distinct highs additionally
permits selected earlier highs in old linear inputs, and permits its
OR node to be selected as a high. Pure-product classes use activity
bits: an unselected prior product may be read alone, but mixing it
with another argument requires its activity bit. Context and pair
selectors imply the relevant activity bits. A relaxation bit permits
the one exceptional merged context in the repeated-high class.
All classes allow empty clauses, zero or constant-one contexts,
repeated source choices, unused nodes and restricted coincidences.

The seven exact refutations have the following sizes, in the order
of the classes specified above:
\[
\begin{array}{c|l|r|r|r}
 &\text{class}&\text{variables}&\text{clauses}&\text{proof lines}\\\hline
1&3\text{ ANDs},\ 2\text{ private ORs},\ 1\text{ cut}&1445&4932&19292\\
2&3\text{ ANDs},\ 1\text{ private OR},\ 2\text{ distinct highs}&1669&5856&19529\\
3&3\text{ ANDs},\ 1\text{ private OR},\ 1\text{ repeated high}&1130&3836&1678\\
4&3\text{ ANDs},\ 3\text{ distinct visible highs}&1595&5658&2912\\
5&3\text{ ANDs},\ 2\text{ distinct highs in }3\text{ cuts}&995&3491&1277\\
6&4\text{ ANDs},\ 1\text{ private OR},\ 1\text{ cut}&1445&4928&5200\\
7&4\text{ ANDs},\ 2\text{ distinct visible highs}&1389&4953&2932.
\end{array}
\]
Each CNF was independently regenerated byte for byte, and every
refutation was checked in RUP-only mode with zero RAT steps. The
checks took between $0.11$ and $0.42$ seconds. The program
\emph{Six-deletion path obstruction}, exact source bindings,
compressed refutations and checker logs accompany the workbench.
Together with the two-product and one-addition exclusions, these
classes exhaust the six-deletion cases. The last assertion follows
because the only possible counterexamples with an optimal frame and
cover gain three had five or six deletions.
\end{proof}
